\documentclass{article}
\usepackage{graphicx} % Required for inserting images
\usepackage{amsmath}
\usepackage{amssymb}
\usepackage{xcolor}
\usepackage{float}
\usepackage{listings}
\usepackage{xcolor}

\lstset{
    % change caption to be below the listing
    % choose appropriate font style.
    captionpos=b,
  basicstyle=\footnotesize\ttfamily,
  numbers=left,
  numberstyle=\tiny,
  stepnumber=1,
  numbersep=8pt,
  frame=single,
  breaklines=true,
  tabsize=2,
  showstringspaces=false,
  keywordstyle=\color{blue},
  commentstyle=\color{gray},
  stringstyle=\color{teal},
}

\title{Pakula \& Sigel (1970) - Summary \& Questions}
\author{Tal Ramot}
\date{November 2025}

\begin{document}

\maketitle

\section{Introduction}
\subsection{Statement of the problem}
\subsubsection{A The hydrodynamic phenomenon of a shock wave}
The process described in the introduction goes like:
\begin{itemize}
    \item A supersonic heat wave propagates through a solid material with density $\rho_s$.
    \item The velocity of the heat wave decreases as it travels through the material.
    \item A hydrodynamic expansion of the heated material into the vaccum begins.
    \item When the velocity of the front (heat wave) drops below the speed of sound in the material, 
    it is taken by a shock wave (ablative heat wave). 
\end{itemize}
Then, there is a discussion about the conduction approximation. 
\textcolor{red}{understand!}

\subsection{Hydrodynamic Equations with nonlinear heat conduction}
\subsubsection{Definition of Physical attributes \& unit analysis}
The physical attributes used in the equations are:
\begin{itemize}
    \item $x$ - distance from the surface $[cm]$
    \item $t$ - time since the beginning of the process $[s]$
    \item $\rho$ - volume density of the material $[g/cm^3]$
    \item $u$ - velocity of the material gathered mass $[cm/s]$
    \item $p$ - pressure of the material gathered mass $[dyne/cm^2]=[g/(cm\cdot s^2)]$
    \item $e$ - specific internal energy $[erg/g]$
    \item $T$ - temperature of the material $[K]$
    \item $\kappa$ - heat conduction coefficient $[erg/(cm \cdot s \cdot K)]$
\end{itemize}

\subsubsection{The hydrodynamic equations}
\underline{Conservation of \textbf{MASS}}
\begin{equation}
    \frac{\partial v}{\partial t} - \frac{\partial u}{\partial m} = 0
\end{equation}
where $v = \frac{1}{\rho}\quad([u] = [cm/s])$ is the specific volume, i.e. volume per unit mass, 
and $m = \int_0^x \rho dx\quad([m] = [g/cm^2])$ is the Lagrangian mass coordinate describing the gathered mass up to distance $x$.
This means that \textit{a change in the specific volume with time is due to a change in velocity of that gathered mass.}
\\ \\
\underline{Conservation of \textbf{MOMENTUM}}
\begin{equation}
    \frac{\partial u}{\partial t} - \frac{\partial p}{\partial m} = 0
\end{equation}
where $[u] = [cm/s]$ is the velocity of the gathered mass and $[p] = [g/(cm \cdot s^2)]$ is the pressure (Force per unit Area)
that the mass exerts. Notice that $\frac{\partial}{\partial m} = \frac{1}{\rho} \frac{\partial}{\partial x}$
\textit{a change in the velocity of the gathered mass with time is due to a change in pressure of that gathered mass.}
\\ \\ 
\underline{Conservation of \textbf{ENERGY}}
\begin{equation}
    \frac{\partial e}{\partial t} + p \frac{\partial v}{\partial t} = \frac{\partial}{\partial m} \left( \kappa \frac{\partial T}{\partial m} \right)
\end{equation}
where $\frac{\partial}{\partial m} \left( \kappa \frac{\partial T}{\partial m} \right) = -\frac{\partial S}{\partial m}$.
$S$ is the heat flux (energy per unit area per unit time) and $\kappa$ is the heat conduction coefficient. 
\textit{a change in the internal energy of the gathered mass with time plus the work done by pressure due to change in specific volume
is equal to the change in heat flux with respect to the gathered mass. (First law of thermodynamics).}


\subsubsection{THe heat internal energy \& equation of state}
\begin{equation}
    e = \frac{1}{\gamma - 1}pv
\end{equation}
\underline{Note:} This equation is derived from finding the internal energy gained
by an ideal gas when in an adiabatic expansion process.
Important things to use (may ask GPT about): Mayer's relation $C_p - C_v = R$, definition of internal energy $de = C_v dT$,
and the adiabatic process relation $pv^\gamma = const$.
\begin{equation}
    pv = KT^\delta
\end{equation}
and the constants $a, \lambda$ are both defined as the following wall:
\begin{equation}
    a = \frac{\kappa_0}{\delta \kappa^{\frac{\nu + 1}{\delta}}}
\end{equation}
\begin{equation}
    \lambda = \frac{\nu + 1 - \delta}{\delta}
\end{equation}
substituting the above equations for $e$ and $pv$ in the energy conservation equation, we get:
\begin{equation}
    \frac{1}{\gamma - 1}  \frac{\partial (pv)}{\partial t} + p \frac{\partial v}{\partial t} = \frac{\partial}{\partial m} \left( a p^{\nu} v^{\lambda} \frac{\partial T}{\partial m} \right)
\end{equation}
or simplified:
\begin{equation}
    \frac{\partial p}{\partial t} + \gamma p \frac{\partial v}{\partial m} = a(\gamma - 1) \frac{\partial}{\partial m} \left( p^{\nu} v^{\lambda} \frac{\partial T}{\partial m} \right)
\end{equation}

\subsection{Self Similarity - Theory}
\subsubsection{A proof for the Buckingham Pi's Theorem}
\textcolor{red}{To be continued...}

\subsubsection{Explaining the derivations for the self similarity analysis in this subsection}
Consider a physical quantity $g(a_1,...,a_k,a_{k+1},...,a_n) = 0$ where $a_1,...,a_k$ have independent dimensions.
\underline{statement:} $g$ can be rewritten as:
\begin{equation}
    g = g_1 (a_1,...,a_k) \tilde{G}(\xi_1,...,\xi_{n-k}) 
\end{equation}
where $\xi_i$ are dimensionless quantities defined as:
\begin{equation}
    \xi_i = a_{k+i} a_1^{x_1} a_2^{x_2} ... a_k^{x_k}
\end{equation}
and $g_1$ is the dimensional part of $g$ defined by:
\begin{equation}
    g_1 = a_1^{\alpha_1} a_2^{\alpha_2} ... a_k^{\alpha_k}
\end{equation}
\underline{proof outline:} Suppose a = $(a_1,...,a_n)$ is the units (dependent) of the physical problem, and that
$D\in R^{k \times n}$ is the dimension matrix of the problem. Then, changing units corresponds to the transformation:
\begin{equation}
    log(a) \to log(a) + D^T \lambda, \quad \lambda \in R^k
\end{equation}
which is a group action of $R^k$ on $R^n$. Let $g(a) = 0$ be a physical law. It must be invariant under this group action, i.e.:
\begin{equation}
    g(a) = g(a') \quad \text{where} \quad log(a') = log(a) + D^T \lambda
\end{equation}
and we must remember that the transformation of g is:
\begin{equation}
    g(a') = g(a) + \sum_{i=1}^n \frac{\partial g}{\partial log(a_i)} \sum_{j=1}^k D_{j,i} \lambda_j
\end{equation}
which means that from Linear Algebra we get that:
\begin{equation}
    \mathbb{R}^n = \text{Im}(D^T) \oplus \text{Ker}D
\end{equation}
and therefore, we can choose a basis for $\mathbb{R}^n$ such that the first k basis vectors span Im$(D^T)$ (with dimension) and the last n-k basis vectors span Ker$(D)$ (dimensionless).
which gives us:
\begin{equation}
    g(a) = a_1^{\alpha_1} a_2^{\alpha_2} ... a_k^{\alpha_k} \tilde{G}(\xi_1,...,\xi_{n-k})
\end{equation}

\subsection{Self Similarity of the 2nd kind (Assymptotic)}
Here the text explains that besides the physical variables with dimensions $t, a, W_\nu^*, m$ - there is $\rho_s$.
This extra variable makes the "dimensionless group" (achieved by Buckingham Pi's theorem) to be 2-dimensional.
However, when $t \to \infty$, the influence of $\rho_s$ overwhelms the initial density $\rho_0$ and the remaining solution
is fully self-similar.

\subsection{}

\section{Questions \& Tough Points}
\begin{enumerate}
    \item "The range of validity of the conduction approximation is 
    discussed at the end of the paper (Sec V or VI)" 
    - I suppose that it's referring to the relation of $K \sim T^\nu$,
    I need to make sure of that and read the last paragraph of the introduction.
    \item research on the derivation of the EOS and internal energy - they both might
    come from basic thermodynamics.
    \item How deeply should I understand the self-similarity derivation form Buckingham Pi's theorem shown in subsection 2.C?
    Specifically the derivation of $g, g_1, \xi_i$? Chat GPT said that it involves many deep claims from group theory and linear algebra.
    \item Relates to the subsection of thermodynamics relations - what exactly is the thermodynamical process that is discussed
    there? is it adiabatic expansion? is it something else? How is is related to the given relations for $e$ etc.
    \item The big picture - Undestand the goal of this article - Explain how to take hydrodynamic equations
    and solve it using self-similarity of the 2nd kind (explained in the article), and with the approximations of the front boundary conditions.


\end{enumerate}


\section{Meeting with Shi - Week 12, Monday}
\begin{enumerate}
    \item The article has a corrction called "Euratom", about some of the power rules presented.
    The Erratum still constains mistakes, though. For example, figure 1 describing the temperature profile misincludes
    the flat region of the shockwave at the supersonic regime.
    Despite the Euratom, the article is considered revolutinary as it discusses a solution of the subsonic region.
    \item Relating to the figures - there is a special case in which the shock wave and the temperature wave move 
    in the same speed. Need to read the course summary from Prof. Kreif to fully undertsand.
    \item Shussman wrote solvers of the subsonic, sonic, shockwave and the connection-edges of the shockwave in Matlab -
    after completing the article, \textcolor{red}{ask Shi for the code}.
    \item Phonons - they are modeling the ions in a lattice. In our article, the conduction is performed through
    the interaction between the ions. In Pomraning, the conduction is through photons (emission-absorption, who are 
    encapsuled in the Boltzmann equation in $\sigma_a$, $\sigma_s$ which are determined by QM).
    I need to understand better the difference between the two models.
    \item In a hohlraun, in the incident point of the laser, the plasma is very diluted so the electronic conduction is negligible? 
    understand better.
    \item Understand what is the $\kappa$ term in the equation of heat conduction. Read on appendix C.
    
\end{enumerate}


\end{document}